\documentclass{article}

\begin{document}
    \begin{center}
        \Large \textbf{Challenge 3: Terraform CD Pipeline} \\ [6pt]
        \normalsize
        Autor: Maximilian Jakob Maag B.Sc. \\
        Version: 1.0
    \end{center}
    \section{Intro}
    Terraform in einer Enterprise Umgebung muss  in der Lage sein von mehereren Personen gleichzeitig verwendet zu werden.
    Eine Lösung für dieses Problem ist der Einsatz von Git. Ein Git Repository speichert den Terraform Code. Die einzelnen Admin/Entwickler reichen Pull Requests ein, um Änderungen zunächst in 
    einem Dev/Staging Branch zusammen zu führen und zu testen. Dieser wird anschließend in den Main Branch gemerged, was eine CD Pipeline auslöst, welche terraform apply auf die Produktivumgegung anwendet.
    CI/CD ist eine wichtige Praxis aus der Softwareentwicklung und in modernen Setups nahezu unverzichtbar.
    \\ \\
    Für diese Challenge ist die Zusammenarbeit mit mehreren Personen erforderlich. Eine Person errichtet die CD Pipeline.
    Die Anderen Personen schlagen Änderungen mittels PR vor. Nach erfolgreichem Merge mus die Infrastruktur von der Pipeline deployed werden.
    \section{Aufgaben}
    Diese Challenge soll Ihnen dabei helfen Terraform im Enterprise Umfeld einsetzen zu können.
    Dafür errichten Sie mehrere CD Pipelines für eine Staging und eine Produktivumgebung.

    \subsection{Aufgabe 1}
    Legen Sie ein Git Repository an. Dieses Repository muss auf einem Git Server liegen.  \\ \\
    Errichten Sie eine multi environment Umgebung in Terraform ein. Es muss mindestens eine Stageing und eine Prod Umgebung vorhanden sein.

    \subsection{Aufgabe 2}
    Erstellen Sie mindestens ein Modul, welches in beiden Umgebungen verwendet werid. Ein simpler Docker Stack ist ausreichend.

    \subsection{Aufgabe 3}
    Sobald ein Pull Request in den Stageing Branch gemerged werden soll.  Soll eine CI Pipeline überprüfen, ob der Terraform Code korrekt kompiliert.
    Nach dem der Merge erfolgt ist soll eine CD Pipelien die Infrastruktur deployen.

    \subsection{Aufgabe 4}
    Setzen Sie Aufgabe 3 in der Prod Umgebung um. Sobald ein Pull Request in den Main Branch gemerged wird, soll die Infrastruktur in der Prod Umgebung deployed werden.
    Es soll nur von einem Staging Branch in den Main Branch gemerged werden. Das pushen von Änderungen direkt auf den Main Branch soll unterbunden werden.

    \subsection{Aufgabe 5}
    Testen Sie Ihre CD Pipeline, indem Sie eine Änderung an Ihrem Modul vornehmen.
    Gehen Sie einen Infrastruktur Change vollständig durch. Führen Sie zunächst in der Staging Umgebung aus. Anschließend testen Sie dort Ihre Änderung und überführen diese in den Main Branch.
    Die Übung ist erfolgreich, wenn Ihnen diese Änderung ohne Desaster oder Rollback gelungen ist. \\ \\
    Nach dem Abschluss der Übung führen Sie in allen Umgebungen terraform destroy aus.
\end{document}